\chapter{Antipattern nelle architetture a microservizi}



\section{Antipatterns e code smells}
Come già discusso nelle sezioni precedenti, i microservizi stanno guadagnando sempre più popolarità e importanza grazie alle loro caratteristiche. Per facilitare tale sviluppo, come riportato nel paragrafo~\ref{sec:par2_4}, sono state identificate delle soluzioni a problemi di progettazione ricorrenti, tali soluzione prendono il nome di \textit{design patterns} \cite{CernyCatalogo}.

Analogamente, gli \textbf{anti-patterns} sono identificabili come pratiche o scelte di progettazione a problemi ricorrenti che, pur sembrano efficienti in una prima fase, producono successivamente conseguenze negative e compromettono di conseguenza la qualità complessiva del sistema \cite{CernyCatalogo}.

Potremmo quindi definirli, come soluzioni \textit{"non ottimali"} a problemi ricorrenti. Un bad smell (o code smell), invece, rappresenta un sospetto che qualcosa non sia corretto e si configura a tutti gli effetti come un segnale di allarme che fa trasparire eventuali problemi nel design, ma non fornisce soluzione specifica \cite{CernyCatalogo}.
Citando Fowler, un code smells è "come un segnale superficiale che spesso indica un problema più profondo nel sistema" \cite{fowler_codesmell}.
Potremmo quindi identificare gli anti-pattern come la causa profonda del problema mentre i bad smells indicano i sintomi osservabili. 

L'origine degli anti-pattern è dovuta a varie cause che non sono solo vincolate ad aspetti tecnici ma anche a situazioni culturali e puramente organizzative:
\begin{itemize}
    \item Pressioni dovute a scadenze del progetto che portano a decisioni frettolose per concludere entro i termini previsti. 
    \item Mancanza di esperienza e competenza da parte del team.
    \item Mancanza di comunicazione tra i membri del team o tra stakeholder e sviluppatori. 
    \item Requisiti poco chiari, incompleti o in continua evoluzione. 
    \item Eccessiva fiducia in soluzioni adottate precedentemente con conseguente riutilizzo, ma come detto in precedenza le soluzioni vanno adottate in base al contesto di sviluppo che varia a seconda dei casi. 
    \item Scarsa documentazione del sistema. 
    \item Mancanza di test o testing inadeguato. 
\end{itemize}

Gli anti-pattern hanno ovviamente un impatto sul sistema. Le 

conseguenze che possiamo riscontrare sono le seguenti:
\begin{itemize}
    \item Aumento dell'accoppiamento tra i servizi, 
    
    compromettendo l'indipendenza.
    \item Aumento della complessità complessiva del sistema. 
    \item Difficoltà nel mantenere e nel testare il sistema, dovuto alla presenza di dipendenze implicite e interazioni non gestite opportunamente. 
    \item Degrado delle prestazioni, spesso causato da comunicazioni eccessive tra microservizi o da scelte architetturali inefficienti.
    \item Perdita di scalabilità a causa dei servizi che diventano dei veri e propri colli di bottiglia  difficili da isolare e ottimizzare. 
\end{itemize}

Ad esempio, nel caso di cyclic dependency, due microservizi finiscono per dipendere reciprocamente e il fallimento di uno si propaga inevitabilmente all’altro. Analogamente, con il shared database più microservizi accedono allo stesso database, aumentando l’accoppiamento e riducendo l’autonomia dei servizi.

In questo elaborato si fa riferimento alla convenzione terminologica adottata da Tomas Cerny e Davide Taibi, nella quale il termine anti-pattern è impiegato come categoria ombrello comprendente anche i bad smells, per ragioni di semplificazione concettuale \cite{CernyCatalogo}.

\section{Catalogo degli anti-pattern}

Lo studio condotto da Cerny e Taibi~\cite{CernyCatalogo} ha permesso di definire un catalogo composto da 58 anti-pattern, suddivisi in cinque macro-categorie:
\begin{itemize}
  \item \textbf{Intra-service design};
  \item \textbf{Inter-service decomposition};
  \item \textbf{Service interaction};
  \item \textbf{Security};
  \item \textbf{Team organization}.
\end{itemize}


La Tabella~\ref{tab:catalogo_antipattern} sintetizza gli anti-pattern appartenenti alle cinque macro-categorie del catalogo, riportando per ciascuno una breve descrizione.
Sono analizzate nel dettaglio nelle sezioni successive le singole categorie.


\subsection{Intra-service design}

Gli anti-pattern appartenenti a questa categoria fanno riferimento a difetti di progettazione interna di un singolo servizio. In particolare, riguardano aspetti come la granularità del servizio, la qualità dell’interfaccia esposta e la coesione funzionale delle responsabilità implementate.

La categoria \textit{Intra-service design} è suddivisa, in generale, in tre sottocategorie: \textit{granularity}, \textit{service interface} e \textit{functional cohesion}~\cite{CernyCatalogo}.

\begin{itemize}
    \item \textit{Granularity}: riguarda la dimensione del servizio, includendo casi come \textit{Nano-service} e \textit{Mega-service}.
    \item \textit{Service interface}: riguarda problemi relativi all’interfaccia del microservizio, come \textit{CRUDY service}, \textit{Nobody home}, \textit{Data service} e \textit{No API-versioning}.
    \item \textit{Functional cohesion}: riguarda la chiarezza, la comprensibilità e la coesione funzionale del servizio, includendo casi come \textit{Low cohesive operation} e \textit{Ambiguous service}.
\end{itemize}

Appartengono alla sottocategoria \textit{granularity} i seguenti anti-pattern:
\begin{itemize}
    \item \textbf{Nano-service}: microservizi frammentati eccessivamente, al punto che il sovraccarico di comunicazione di rete e manutenzione supera i benefici derivanti dalla loro separazione.
    \item \textbf{Mega-service}: servizio di dimensioni eccessive, con troppe responsabilità e una quantità di codice tale da compromettere scalabilità, indipendenza e manutenibilità.
\end{itemize}

Appartengono alla sottocategoria \textit{service interface} i seguenti anti-pattern:
\begin{itemize}
    \item \textbf{CRUDY service}: servizio che espone principalmente operazioni CRUD, senza una reale logica di business.
    \item \textbf{Nobody home}: servizio deployato e funzionante, ma non utilizzato o invocato da altri componenti del sistema. Si tratta quindi di un servizio che richiede risorse in termini di manutenzione, ma che non porta alcuna utilità effettiva.
    \item \textbf{Data service}: servizio molto semplice che espone principalmente metodi per l’accesso ai dati, spesso tramite parametri primitivi.
    \item \textbf{No API-versioning}: mancanza di versionamento delle API, con il rischio di introdurre breaking changes per i client esistenti.
\end{itemize}

Appartengono alla sottocategoria \textit{functional cohesion} i seguenti anti-pattern:
\begin{itemize}
    \item \textbf{Whatever types}: uso di tipi di dati generici che rendono ambiguo il contratto dell’interfaccia.
    \item \textbf{Low cohesive operation}: presenza di operazioni con bassa coesione raggruppate nello stesso servizio.
    \item \textbf{Ambiguous service}: servizio con responsabilità non ben definite, che può causare confusione e sovrapposizione di funzionalità.
\end{itemize}

\subsection{Inter-service decomposition}

Gli anti-pattern appartenenti a questa categoria riguardano errori nella decomposizione del sistema in servizi e nelle relazioni strutturali tra essi. Nel dettaglio si basano su scelte errate nella suddivisione del sistema in microservizi, accoppiamenti stretti o violazione dei principi di modularità che caratterizzano architetture di questo genere.

La categoria \textit{Inter-service decomposition} include quindi problematiche che emergono non tanto all’interno del singolo servizio, quanto nel modo in cui i servizi vengono separati, organizzati e integrati tra loro~\cite{CernyCatalogo}.

Appartengono a questa categoria i seguenti anti-pattern:
\begin{itemize}
  \item \textbf{Transactional integration}: transazioni che si estendono oltre i confini del singolo servizio, violandone l’autonomia.
  \item \textbf{Co-change coupling}: servizi che devono essere modificati frequentemente insieme a causa di un forte accoppiamento.
  \item \textbf{Duplicated services}: servizi multipli che implementano funzionalità simili o identiche.
  \item \textbf{Microservice greedy}: creazione eccessiva di microservizi senza chiari vantaggi misurabili.
  \item \textbf{Service chain}: lunga catena di chiamate tra servizi, con conseguente aumento di latenza e complessità.
  \item \textbf{Hub-like dependency}: servizio che diventa un punto centrale di dipendenza per molti altri servizi.
  \item \textbf{Cyclic dependency}: dipendenze circolari tra servizi, che creano problemi di deployment e manutenzione.
  \item \textbf{Chatty service}: comunicazione eccessiva e troppo fine-grained tra servizi, con effetti negativi sulle prestazioni.
  \item \textbf{Shared persistency}: più servizi che accedono allo stesso database, generando accoppiamento a livello dei dati.
  \item \textbf{Sand pile}: sistema composto da molti nano-service che condividono dati comuni.
  \item \textbf{Shared libraries}: condivisione di librerie runtime tra microservizi, con il rischio di comprometterne l’indipendenza.
  \item \textbf{Wrong cuts}: decomposizione basata su layer tecnici invece che sul dominio di business.
  \item \textbf{Knot service}: servizi fortemente accoppiati tramite integrazioni point-to-point.
  \item \textbf{Scattered parasitic functionality}: funzionalità distribuite su più servizi, con responsabilità sovrapposte.
\end{itemize}

\subsection{Service interaction}

Gli anti-pattern appartenenti a questa categoria riguardano problemi nei meccanismi di comunicazione tra i servizi. Tale categoria di problemi può compromettere fattori importanti per il sistema come: la scalabilità, la manutenibilità e la capacità di evoluzione.

La categoria \textit{Service interaction} include quindi difetti legati al modo in cui i servizi si invocano, scambiano messaggi, gestiscono endpoint, fallimenti e informazioni operative~\cite{CernyCatalogo}.

Gli anti-pattern associati a questa categoria sono:
\begin{itemize}
  \item \textbf{ESB misuse}: uso di un Enterprise Service Bus come punto centrale di integrazione, con il rischio di creare un bottleneck.
  \item \textbf{On-line only}: sistema che evita completamente l’elaborazione batch, anche quando questa sarebbe appropriata.
  \item \textbf{Empty messages}: scambio di messaggi vuoti, utilizzati come semplici segnali.
  \item \textbf{Use of business logic in communication}: logica di business incorporata nello strato di comunicazione.
  \item \textbf{Hardcoded endpoint}: indirizzi ed endpoint hardcodati nel codice, invece dell’utilizzo di meccanismi come il service discovery.
  \item \textbf{No API-gateway}: mancanza di un gateway unificato, che costringe i client a chiamare direttamente i singoli servizi.
  \item \textbf{Wobbly service interactions}: interazioni tra servizi che non gestiscono adeguatamente i fallimenti, riducendo la resilienza del sistema.
  \item \textbf{Timeout}: timeout fissi che non considerano adeguatamente i diversi contesti operativi.
  \item \textbf{No health check}: mancanza di endpoint o meccanismi per verificare lo stato di salute dei servizi.
\end{itemize}

\subsection{Security}

Gli anti-pattern di questa categoria riguardano principalmente problematiche inerenti alla sicurezza dei microservizi.

La categoria \textit{Security} riguarda, quindi, la protezione dei sistemi attraverso tre meccanismi fondamentali: autenticazione, autorizzazione e cifratura. L’autenticazione è il processo di verifica dell’identità di un utente; l’autorizzazione determina quali risorse o operazioni sono accessibili a un utente autenticato; la cifratura consente invece di trasformare un testo leggibile in un formato illeggibile, al fine di proteggerne il contenuto~\cite{CernyCatalogo}.

Gli anti-pattern associati a questa categoria sono:
\begin{itemize}
  \item \textbf{Unauthenticated traffic}: traffico API non autenticato tra servizi o da sistemi esterni.
  \item \textbf{Multiple user authentication}: presenza di molteplici punti di accesso per l’autenticazione utente, con conseguente aumento della superficie d’attacco.
  \item \textbf{Publicly accessible microservices}: microservizi direttamente accessibili da client esterni senza adeguata protezione.
  \item \textbf{Unnecessary privileges}: privilegi eccessivi concessi ai microservizi rispetto a quanto effettivamente necessario.
  \item \textbf{Insufficient access control}: controlli di accesso mancanti o inadeguati sui microservizi.
  \item \textbf{Centralized authorization}: autorizzazione centralizzata, invece di una gestione più granulare a livello di singolo servizio.
  \item \textbf{Non-secured service communications}: comunicazioni tra servizi non cifrate o non autenticate.
  \item \textbf{Non-encrypted data exposure}: dati sensibili esposti in chiaro nello storage o durante il trasporto.
  \item \textbf{Own crypto code}: implementazione personalizzata di algoritmi crittografici, invece dell’utilizzo di librerie standard.
  \item \textbf{Hardcoded secrets}: segreti, come password o API key, hardcodati nel codice o nelle configurazioni.
\end{itemize}

\subsection{Team organization}

Gli anti-pattern appartenenti a questa categoria riguardano problemi organizzativi e operativi legati al modo in cui i team progettano, sviluppano, rilasciano e gestiscono i microservizi. A differenza delle categorie precedenti, il focus non è solo sulla struttura tecnica del software, ma anche sulle pratiche adottate durante il ciclo di vita del sistema.

La categoria \textit{Team organization} include quindi difetti legati sia alle scelte dei team di sviluppo sia alla gestione operativa del sistema, comprese eventuali violazioni delle pratiche DevOps~\cite{CernyCatalogo}.

Gli anti-pattern associati alla sottocategoria \textit{development} sono:
\begin{itemize}
  \item \textbf{Shiny nickel}: adozione di nuove tecnologie solo perché di tendenza, senza reali benefici.
  \item \textbf{Golden hammer}: uso della stessa tecnologia familiare per ogni problema, anche quando non appropriata.
  \item \textbf{Lack of communication standards}: assenza di standard comuni per API e formati di messaggio tra team.
  \item \textbf{Too many standards}: eccessiva eterogeneità di linguaggi, protocolli e framework nell’organizzazione.
  \item \textbf{Inadequate techniques support}: uso di tecniche inadeguate, che non supportano efficacemente lo sviluppo di microservizi.
  \item \textbf{Single layer team}: team organizzati per layer tecnici invece che per servizio o dominio.
  \item \textbf{No legacy}: tendenza a eliminare completamente i sistemi legacy, invece di integrarli progressivamente.
  \item \textbf{Data-driven migration}: migrazione simultanea di dati e funzionalità durante la scomposizione del sistema.
  \item \textbf{Big bang}: riscrittura completa del sistema in una sola volta, invece che tramite un processo incrementale.
  \item \textbf{Multiple service instances per host}: deployment di più istanze di servizi sullo stesso host, con conseguente condivisione delle risorse.
\end{itemize}

Gli anti-pattern associati alla sottocategoria \textit{operation} sono:
\begin{itemize}
    \item \textbf{No CI/CD}: mancanza di integrazione e deployment continui, con processi prevalentemente manuali.
    \item \textbf{Manual configuration}: configurazione manuale invece di automazione e strumenti di configuration management.
    \item \textbf{Insufficient monitoring}: monitoraggio inadeguato delle prestazioni e dei failure dei servizi.
    \item \textbf{Dismiss documentation}: documentazione insufficiente delle API e dell’architettura.
    \item \textbf{Insufficient message traceability}: messaggi privi di metadati sufficienti per tracciare origine e percorso.
    \item \textbf{Local logging}: log scritti solo in storage locale invece che inviati a sistemi centralizzati.
\end{itemize}

